% Solutions of Book 4, chapter 1 (cpp-classes-raii).
\section*{Chapter \ref{ch:b2:cpp-classes-raii} --- From C to C++: Classes and Resource Management}

\begin{solution}{exo:b2:cpp-classes-raii:1}
The global \omterm{def:b2:cpp-classes-raii:class}{object} is built before \code{main} and destroyed after it; inside, \code{f(1)} makes \code{x1}, calls
\code{f(0)}, which makes and drops \code{x0}, then drops \code{x1}; the inner block makes and drops \code{b}:
\omoutput{code/bachelor-2/01-cpp-classes-raii/outputs/exo_order.txt}{The ten lines.}
\end{solution}

\begin{solution}{exo:b2:cpp-classes-raii:2}
\code{result} is an automatic \omterm{def:b2:cpp-classes-raii:class}{object} of \code{longer}: it dies when \code{longer} returns, so the returned \omterm{def:b2:cpp-classes-raii:reference}{reference}
names an \omterm{def:b2:cpp-classes-raii:class}{object} that no longer exists (a \emph{dangling} \omterm{def:b2:cpp-classes-raii:reference}{reference}), and reading it in \code{main} is undefined
behaviour. Any outcome is allowed; this compiler, knowing the \omterm{def:b2:cpp-classes-raii:reference}{reference} is dangling, returned a null address and the
print produced nothing. The simplest fix returns a copy: declare the return type \code{std::string} instead of
\code{const std::string &}, and the body can shrink to \code{return a.size() >= b.size() ? a : b;}. (Returning \code{a} or \code{b} by \omterm{def:b2:cpp-classes-raii:reference}{reference} would also be
correct, since they name the caller's \omterm{def:b2:cpp-classes-raii:class}{objects}, but only as long as those outlive the use.)
\end{solution}

\begin{solution}{exo:b2:cpp-classes-raii:3}
\omcode{code/bachelor-2/01-cpp-classes-raii/exercises/int_array.hpp}{7}{27}{\code{IntArray}}
Step 1: the \omterm{def:b2:cpp-classes-raii:constructor}{constructor} acquires the block (\code{new int[n]{}} zero-fills it, and throws \code{std::bad_alloc} if
there is no memory, so no \omterm{def:b2:cpp-classes-raii:class}{object} without a block exists). Step 2: the \omterm{def:b2:cpp-classes-raii:constructor}{destructor}'s \code{delete[]} is the only
release. Step 3: the two deleted functions. Step 4: \code{size} and the two \code{at}, which keep every access inside
the block. Tested in \code{tests/exercises_test.cpp}.
\end{solution}

\begin{solution}{exo:b2:cpp-classes-raii:4}
Each of the $nm$ \omterm{def:b2:cpp-classes-raii:class}{objects} prints one \code{make} and one \code{drop} line: $2nm$ lines. When \code{n.m} is made, all
$nm$ \omterm{def:b2:cpp-classes-raii:class}{objects} are alive (the blocks are nested and none has ended), so by \cref{prop:b2:cpp-classes-raii:reverse-order}
they die in the reverse order of construction. Construction order is \code{1.1}, \dots, \code{1.m}, \code{2.1},
\dots, \code{n.m}, so the first drops are \code{n.m}, \code{n.(m-1)}, \code{n.(m-2)} and the last three
\code{1.3}, \code{1.2}, \code{1.1} (for $m \geq 3$).
\end{solution}

\begin{solution}{exo:b2:cpp-classes-raii:5}
\omoutput{code/bachelor-2/01-cpp-classes-raii/outputs/exo_handlers.txt}{The output.}
$k = 0$: nothing thrown. $k = 1$: \code{invalid_argument} skips the \code{out_of_range} handler (not its type) and is
caught by the second inner handler, which prints and rethrows with \code{throw;}; the middle handler catches it,
since \code{invalid_argument} derives from \code{logic_error}, and \code{step done} is never reached. $k = 2$: the
first inner handler matches; the \omterm{def:b2:cpp-classes-raii:exception}{exception} is finished, so \code{step} goes on and prints \code{step done}.
$k = 3$: \code{runtime_error} matches neither inner handler nor \code{logic_error}; the outer
\code{std::exception} handler catches it. $k = 4$: an \code{int} is not a \omterm{def:b2:cpp-classes-raii:class}{class} derived from \code{std::exception};
only \code{catch (...)} matches.
\end{solution}

\begin{solution}{exo:b2:cpp-classes-raii:6}
Each time, the first \code{fopen} succeeds, the second fails, and the \omterm{def:b2:cpp-classes-raii:constructor}{constructor} throws. A \omterm{def:b2:cpp-classes-raii:constructor}{constructor} that throws
never produces an \omterm{def:b2:cpp-classes-raii:class}{object}, so the \omterm{def:b2:cpp-classes-raii:constructor}{destructor} never runs, and the raw pointer \code{a_} has no \omterm{def:b2:cpp-classes-raii:constructor}{destructor} of its own:
one file leaks per failure, 4, 5, \dots, 8 open files. The fix makes each file a member that owns itself:
\omcode{code/bachelor-2/01-cpp-classes-raii/exercises/merger.hpp}{5}{14}{\code{Merger} with \code{File} members}
When \code{b_}'s \omterm{def:b2:cpp-classes-raii:constructor}{constructor} throws, the members already built (here \code{a_}) are destroyed before the \omterm{def:b2:cpp-classes-raii:exception}{exception}
leaves the \omterm{def:b2:cpp-classes-raii:constructor}{constructor}, so \code{a_} closes its file. The test opens and fails ten times and finds the count of open
files unchanged.
\end{solution}

\begin{solution}{exo:b2:cpp-classes-raii:7}
\omcode{code/bachelor-2/01-cpp-classes-raii/exercises/scope_guard.hpp}{5}{20}{\code{ScopeGuard}}
In a C-style function, right after each acquisition succeeds, declare its guard:
\code{char *a = (char *)malloc(n); if (!a) return -1; ScopeGuard ga([&] {free(a);});}, then the same for the second
block and for the file. Every later \code{return} (or \omterm{def:b2:cpp-classes-raii:exception}{exception}) runs the guards declared so far, newest first, so
each path frees exactly what was acquired. When the function succeeds and returns the \omterm{def:b2:cpp-classes-raii:raii}{resources} to its caller, they
must not be freed: the function calls \code{dismiss()} on each guard just before the final \code{return}.
\end{solution}

\begin{solution}{exo:b2:cpp-classes-raii:8}
Each of the $k$ checks has its \code{return}, and the end of the body is one more: $k + 1$ paths. In C, path $i$ must
release the $a_i$ \omterm{def:b2:cpp-classes-raii:raii}{resources} held at check $i$, and the last path all $r$: the author writes
$a_1 + \dots + a_k + r$ release calls. With one RAII \omterm{def:b2:cpp-classes-raii:class}{object} per \omterm{def:b2:cpp-classes-raii:raii}{resource}, the function contains none: the release is
written once, in each \omterm{def:b2:cpp-classes-raii:class}{class}'s \omterm{def:b2:cpp-classes-raii:constructor}{destructor}. By \cref{prop:b2:cpp-classes-raii:every-path}, on path $i$ exactly the
$a_i$ \omterm{def:b2:cpp-classes-raii:class}{objects} built so far are destroyed, each once, which is exactly the set the C author had to release by hand.
\end{solution}

\begin{solution}{exo:b2:cpp-classes-raii:9}
\omcode{code/bachelor-2/01-cpp-classes-raii/exercises/bounded_stack.hpp}{7}{17}{\code{BoundedStack}, up to
\code{push}}
The full check comes before any change. The \omterm{def:b2:cpp-classes-raii:constructor}{constructor} reserves room for $N$ elements, so \code{push_back} never
reallocates: it copies \code{x} into the free slot after the last element, and only then increases the size. If the
copy throws, the size was not changed and no existing element was touched, so the stack is exactly as before: the
strong guarantee. (A vector that had to grow would copy the old elements into a new block first, which is where a
weaker implementation could lose them.) \code{pop} and \code{top} throw \code{std::out_of_range} on an empty stack.
\end{solution}

\begin{solution}{exo:b2:cpp-classes-raii:10}
\omcode{code/bachelor-2/01-cpp-classes-raii/exercises/txn_file.hpp}{8}{44}{\code{TxnFile}}
A function that writes its report through a \code{TxnFile} and commits at its end gives the strong guarantee for
the file \code{path}: either the commit happened and the whole report is there, or the \omterm{def:b2:cpp-classes-raii:constructor}{destructor} removed the
temporary and \code{path} is as before. Between the close and the rename, \code{rename} itself may fail (no
permission, or a temporary on another file system); then \code{commit} throws, the \omterm{def:b2:cpp-classes-raii:constructor}{destructor} removes the temporary,
and \code{path} is untouched. A crash of the whole machine there leaves a stale \code{path.tmp}, but never a half
file at \code{path}. The \omterm{def:b2:cpp-classes-raii:constructor}{destructor} runs during unwinding, when another \omterm{def:b2:cpp-classes-raii:exception}{exception} is in flight; letting a second one
out would end the program (\omterm{def:b2:cpp-classes-raii:constructor}{destructors} are \code{noexcept}), so it ignores the results of \code{fclose} and
\code{remove}.
\end{solution}

\begin{solution}{exo:b2:cpp-classes-raii:11}
$o_k$'s \omterm{def:b2:cpp-classes-raii:constructor}{constructor} did not finish, so $o_k$ never existed and has no \omterm{def:b2:cpp-classes-raii:constructor}{destructor} call (whatever part of $R_k$ it had
acquired, it must release itself, for instance with members, as in \cref{exo:b2:cpp-classes-raii:6}). The \omterm{def:b2:cpp-classes-raii:exception}{exception}
leaves blocks $k, k-1, \dots, 1$ in that order, and each destroys its \omterm{def:b2:cpp-classes-raii:class}{object}: $o_{k-1}, \dots, o_1$, so
$R_{k-1}, \dots, R_1$ are released, in that order, by $k - 1$ \omterm{def:b2:cpp-classes-raii:constructor}{destructor} calls. $R_{k+1}, \dots, R_n$ were never
acquired. If the \omterm{def:b2:cpp-classes-raii:exception}{exception} comes after $o_n$ is built, all $n$ \omterm{def:b2:cpp-classes-raii:class}{objects} are alive, and unwinding destroys
$o_n, \dots, o_1$: $n$ calls, releasing $R_n$ first and $R_1$ last.
\end{solution}

\begin{solution}{exo:b2:cpp-classes-raii:12}
\omcodes{code/bachelor-2/01-cpp-classes-raii/exercises/matrix.hpp}{10-21,48-63}{\code{Matrix}: construction,
assignment and the checks}
The \omterm{def:b2:cpp-classes-raii:constructor}{constructor} establishes the invariant: \code{checked} throws before \code{data_} is built when a size is 0, and
otherwise \code{data_} gets exactly $rc$ elements. The three members are private, so only \omterm{def:b2:cpp-classes-raii:class}{member functions} can
change them. Of those, \code{operator()} changes an element, never a size; \code{transposed} and \code{operator*}
create new matrices through the \omterm{def:b2:cpp-classes-raii:constructor}{constructor} and do not change \code{*this}. The copy \omterm{def:b2:cpp-classes-raii:constructor}{constructor} copies a valid
matrix member by member. The one danger is assignment: the implicit one would assign \code{rows_}, \code{cols_},
then \code{data_}, and if copying \code{data_} ran out of memory, the new sizes would sit with the old data. The
\omterm{def:b2:cpp-classes-raii:class}{class} therefore takes its argument by value (the copy is made first, and may throw before \code{*this} is touched)
and swaps the three members, which cannot throw. No public function can break the invariant. Tested in
\code{tests/exercises_test.cpp}.
\end{solution}

\begin{solution}{pb:b2:cpp-classes-raii:1}
\textbf{1.} Five: the four early returns $-1$ to $-4$ and the final \code{return n}, that is $k + 1$ with $k = 4$.

\textbf{2.} $-1$: nothing is held (\code{fopen} failed). $-2$: \code{f} is held and not closed. $-3$ and $-4$:
\code{f} and \code{buf} are held, neither released. \code{return n}: both are released.

\textbf{3.} 3: the open failed, nothing leaked. 4: the $-99.0$ case returned $-3$ and leaked its \code{FILE} for
\code{/dev/null}. 4: the good batch released everything. 5: \code{/dev/full} returned $-4$ and leaked one more.

\textbf{4.} Each \code{malloc} asked for $16 \cdot 3 + 1 = 49$ bytes, lost in cases 2 and 4: 98 bytes (without
counting the memory of the two leaked \code{FILE}s).

\textbf{5.} Keep the result in a variable and jump to labels placed in the reverse order of acquisition:
\code{free_buf:} (\code{free(buf);}) before \code{close_file:} (\code{fclose(f);}), then one \code{return result;}.
Check 2 jumps to \code{close_file}, checks 3 and 4 to \code{free_buf}: two labels.

\textbf{6.} $r = 2$, $k = 4$, $(a_1, a_2, a_3, a_4) = (0, 1, 2, 2)$: $0 + 1 + 2 + 2 + 2 = 7$ release calls were
needed; the author wrote 2, both on the final path.

\textbf{7.} \code{buf}, a \code{std::string}, owns its heap memory; \code{f}, a \code{File}, owns the \code{FILE}
and the open file. (\code{cell} is an array in the frame: no \omterm{def:b2:cpp-classes-raii:raii}{resource}.)

\textbf{8.} In case 1 the \code{File} \omterm{def:b2:cpp-classes-raii:constructor}{constructor} threw because \code{fopen} failed, so nothing was opened. In case 2
the function threw before \code{f} existed. In case 3 \code{f} died at the end of the function. In case 4 the throw
from \code{flush} unwound the function and destroyed \code{f}, which closed the file. No path leaves a file open.

\textbf{9.} It would have written \code{18.5} before meeting $-99.0$: a half batch at the end of the log, with no
sign that it is incomplete. Checking everything first and writing once leaves the log as it was.

\textbf{10.} The end of the body (line 23); the \code{throw} of line 15; the \code{File} \omterm{def:b2:cpp-classes-raii:constructor}{constructor}'s throw at line
20; \code{write} at line 21; \code{flush} at line 22. A \code{std::bad_alloc} from \code{buf += cell} (line 18) is
possible too.

\textbf{11.} None. The releases are written once, in \code{~File} (its \code{fclose}) and in the \omterm{def:b2:cpp-classes-raii:constructor}{destructor} of
\code{std::string}.

\textbf{12.} \code{f}, built after \code{buf}, is destroyed first and closes the file; then \code{buf} frees its
memory; then the handler in \code{main} runs.

\textbf{13.} The strong guarantee: the function threw before opening the file, so the log and everything else are as
before the call.

\textbf{14.} Only the basic one: nothing leaks and the \code{File} is closed, but on a disk that fills during the
write, part of the batch may already be in the log. (\code{/dev/full} accepts nothing, which hides this.)

\textbf{15.} A \omterm{def:b2:cpp-classes-raii:constructor}{destructor} is \code{noexcept}: an \omterm{def:b2:cpp-classes-raii:exception}{exception} leaving it calls \code{std::terminate}, and it may be
running during the unwinding of another \omterm{def:b2:cpp-classes-raii:exception}{exception}. So \code{~File} cannot report a failed \code{fclose}. A caller
who must know whether the data reached the file calls \code{flush} while the \omterm{def:b2:cpp-classes-raii:class}{object} is alive, where a throw is
allowed; the close that follows has nothing left to write.

\textbf{16.} Open a \code{TxnFile} on the log's name, copy the current log into it, append the batch, then commit,
which renames the copy over the log. If anything throws before the rename, the temporary is removed and the log is
untouched; after the rename, the log holds the old lines and the whole batch.

\textbf{17.} The plain version writes $b$ bytes; the strong one writes $L + b$.

\textbf{18.} $L = 731 \times 5 = 3655$ and $b = 15$: $(3655 + 15)/15 \approx 245$ times more. Once a day, 3.7 kB is
nothing for a disk: the strong version is worth it.

\textbf{19.} Both success timings are below a nanosecond and differ by noise only: being able to throw costs nothing
measurable when nothing is thrown.

\textbf{20.} A failed return code costs about as much as a success, below a nanosecond; a throw and its catch cost
several hundred nanoseconds. The last line of the output gives the ratio of that run: of the order of a thousand.

\textbf{21.} $365 \times 24 \times 60 = 525\,600$ batches a year, of which 5256 fail: about $5$ milliseconds of
throwing in a year.

\textbf{22.} $3 \times 10^7$ throws at about a microsecond: about 30 seconds, against a few hundredths of a second
with error values, a thousand times less.

\textbf{23.} The logger: its failures are rare and handled far from where they occur. The validator should return
an error value per field: an invalid field is a normal outcome there.

\textbf{24.} \emph{A function with $k$ checks has $k + 1$ exit paths by return (more with \omterm{def:b2:cpp-classes-raii:exception}{exceptions}). With every
\omterm{def:b2:cpp-classes-raii:raii}{resource} owned by an \omterm{def:b2:cpp-classes-raii:class}{object}, each acquired \omterm{def:b2:cpp-classes-raii:raii}{resource} is released exactly once on each path, and no release is written
in the function itself: 0 instead of $a_1 + \dots + a_k + r$ (7 for the logger, which wrote 2). On the author's
laptop the success path costs the same with or without \omterm{def:b2:cpp-classes-raii:exception}{exceptions}, below a nanosecond, while a throw costs several
hundred nanoseconds, about a thousand failed return codes.}
\end{solution}
